% Options for packages loaded elsewhere
% Options for packages loaded elsewhere
\PassOptionsToPackage{unicode}{hyperref}
\PassOptionsToPackage{hyphens}{url}
\PassOptionsToPackage{dvipsnames,svgnames,x11names}{xcolor}
%
\documentclass[
  11pt,
]{article}
\usepackage{xcolor}
\usepackage[margin=1in]{geometry}
\usepackage{amsmath,amssymb}
\setcounter{secnumdepth}{-\maxdimen} % remove section numbering
\usepackage{iftex}
\ifPDFTeX
  \usepackage[T1]{fontenc}
  \usepackage[utf8]{inputenc}
  \usepackage{textcomp} % provide euro and other symbols
\else % if luatex or xetex
  \usepackage{unicode-math} % this also loads fontspec
  \defaultfontfeatures{Scale=MatchLowercase}
  \defaultfontfeatures[\rmfamily]{Ligatures=TeX,Scale=1}
\fi
\usepackage{lmodern}
\ifPDFTeX\else
  % xetex/luatex font selection
\fi
% Use upquote if available, for straight quotes in verbatim environments
\IfFileExists{upquote.sty}{\usepackage{upquote}}{}
\IfFileExists{microtype.sty}{% use microtype if available
  \usepackage[]{microtype}
  \UseMicrotypeSet[protrusion]{basicmath} % disable protrusion for tt fonts
}{}
\makeatletter
\@ifundefined{KOMAClassName}{% if non-KOMA class
  \IfFileExists{parskip.sty}{%
    \usepackage{parskip}
  }{% else
    \setlength{\parindent}{0pt}
    \setlength{\parskip}{6pt plus 2pt minus 1pt}}
}{% if KOMA class
  \KOMAoptions{parskip=half}}
\makeatother
% Make \paragraph and \subparagraph free-standing
\makeatletter
\ifx\paragraph\undefined\else
  \let\oldparagraph\paragraph
  \renewcommand{\paragraph}{
    \@ifstar
      \xxxParagraphStar
      \xxxParagraphNoStar
  }
  \newcommand{\xxxParagraphStar}[1]{\oldparagraph*{#1}\mbox{}}
  \newcommand{\xxxParagraphNoStar}[1]{\oldparagraph{#1}\mbox{}}
\fi
\ifx\subparagraph\undefined\else
  \let\oldsubparagraph\subparagraph
  \renewcommand{\subparagraph}{
    \@ifstar
      \xxxSubParagraphStar
      \xxxSubParagraphNoStar
  }
  \newcommand{\xxxSubParagraphStar}[1]{\oldsubparagraph*{#1}\mbox{}}
  \newcommand{\xxxSubParagraphNoStar}[1]{\oldsubparagraph{#1}\mbox{}}
\fi
\makeatother


\usepackage{longtable,booktabs,array}
\usepackage{calc} % for calculating minipage widths
% Correct order of tables after \paragraph or \subparagraph
\usepackage{etoolbox}
\makeatletter
\patchcmd\longtable{\par}{\if@noskipsec\mbox{}\fi\par}{}{}
\makeatother
% Allow footnotes in longtable head/foot
\IfFileExists{footnotehyper.sty}{\usepackage{footnotehyper}}{\usepackage{footnote}}
\makesavenoteenv{longtable}
\usepackage{graphicx}
\makeatletter
\newsavebox\pandoc@box
\newcommand*\pandocbounded[1]{% scales image to fit in text height/width
  \sbox\pandoc@box{#1}%
  \Gscale@div\@tempa{\textheight}{\dimexpr\ht\pandoc@box+\dp\pandoc@box\relax}%
  \Gscale@div\@tempb{\linewidth}{\wd\pandoc@box}%
  \ifdim\@tempb\p@<\@tempa\p@\let\@tempa\@tempb\fi% select the smaller of both
  \ifdim\@tempa\p@<\p@\scalebox{\@tempa}{\usebox\pandoc@box}%
  \else\usebox{\pandoc@box}%
  \fi%
}
% Set default figure placement to htbp
\def\fps@figure{htbp}
\makeatother





\setlength{\emergencystretch}{3em} % prevent overfull lines

\providecommand{\tightlist}{%
  \setlength{\itemsep}{0pt}\setlength{\parskip}{0pt}}



 


\usepackage{amsmath,amssymb,enumitem,booktabs}
\setlist{itemsep=2pt,topsep=3pt}
\usepackage{fancyhdr}\pagestyle{fancy}\fancyhf{}
\fancyhead[L]{\small DS701 Fall 2026}\fancyhead[R]{\small Homework 1}
\fancyfoot[C]{\small\thepage}
\renewcommand{\headrulewidth}{0.4pt}
\usepackage[breakable]{tcolorbox}
% --- answer-space macros (problems PDF only; no-ops in the solutions PDF) ---
% \answerspace{2in}  vertical room for handwritten work, sized from the solution
% \blank             a short fill-in rule; \blank[1.5in] for a wider one
% \answerpage        start the next problem on a fresh page (Gradescope page matching)
\newcommand{\answerspace}[1]{\par\vspace*{#1}}  % \vspace* so room is not dropped at a page break
\newcommand{\blank}[1][1.2in]{\rule{#1}{0.4pt}}
\newcommand{\answerpage}{\clearpage}
\newtcolorbox{solutionbox}{breakable,colback=blue!4,colframe=blue!45!black,title=Solution,fonttitle=\bfseries,left=4pt,right=4pt,top=3pt,bottom=3pt}
\makeatletter
\@ifpackageloaded{caption}{}{\usepackage{caption}}
\AtBeginDocument{%
\ifdefined\contentsname
  \renewcommand*\contentsname{Table of contents}
\else
  \newcommand\contentsname{Table of contents}
\fi
\ifdefined\listfigurename
  \renewcommand*\listfigurename{List of Figures}
\else
  \newcommand\listfigurename{List of Figures}
\fi
\ifdefined\listtablename
  \renewcommand*\listtablename{List of Tables}
\else
  \newcommand\listtablename{List of Tables}
\fi
\ifdefined\figurename
  \renewcommand*\figurename{Figure}
\else
  \newcommand\figurename{Figure}
\fi
\ifdefined\tablename
  \renewcommand*\tablename{Table}
\else
  \newcommand\tablename{Table}
\fi
}
\@ifpackageloaded{float}{}{\usepackage{float}}
\floatstyle{ruled}
\@ifundefined{c@chapter}{\newfloat{codelisting}{h}{lop}}{\newfloat{codelisting}{h}{lop}[chapter]}
\floatname{codelisting}{Listing}
\newcommand*\listoflistings{\listof{codelisting}{List of Listings}}
\makeatother
\makeatletter
\makeatother
\makeatletter
\@ifpackageloaded{caption}{}{\usepackage{caption}}
\@ifpackageloaded{subcaption}{}{\usepackage{subcaption}}
\makeatother
\usepackage{bookmark}
\IfFileExists{xurl.sty}{\usepackage{xurl}}{} % add URL line breaks if available
\urlstyle{same}
\hypersetup{
  pdftitle={DS701 Homework 1: Question Types; Vectors, Norms, Dot Products},
  colorlinks=true,
  linkcolor={blue},
  filecolor={Maroon},
  citecolor={Blue},
  urlcolor={Blue},
  pdfcreator={LaTeX via pandoc}}


\title{DS701 Homework 1: Question Types; Vectors, Norms, Dot Products}
\usepackage{etoolbox}
\makeatletter
\providecommand{\subtitle}[1]{% add subtitle to \maketitle
  \apptocmd{\@title}{\par {\large #1 \par}}{}{}
}
\makeatother
\subtitle{Released Mon Sep 7 · due Wed Sep 16, 11:59 pm on Gradescope}
\author{}
\date{}
\begin{document}
\maketitle


\begin{center}\LARGE\bfseries\color{blue!45!black} SOLUTIONS\end{center}
\renewcommand{\answerspace}[1]{}
\renewcommand{\answerpage}{\par\vspace{6pt}}

\subsection{Problem 1: Classify the
question}\label{problem-1-classify-the-question}

\emph{Practices: naming the question type before choosing a technique
(Lecture 01).}

A music-streaming service is planning its analytics for the fall. Six
questions come up in the planning meeting:

\begin{enumerate}
\def\labelenumi{\arabic{enumi}.}
\tightlist
\item
  What was the median listening time per user last month, broken down by
  country?
\item
  Do our users fall into a few natural groups by the genres they listen
  to --- something worth a closer look?
\item
  Across our whole user base (not just the 2,000 users we surveyed), do
  people who listen during a commute play more podcasts than people who
  do not?
\item
  Will \emph{this} user cancel their subscription in the next 30 days?
\item
  If we give every free-trial user a personalized playlist on day 1,
  will more of them convert to paid?
\item
  By exactly which route does a personalized playlist raise conversion
  --- does it lengthen sessions, which builds a daily habit, which then
  drives payment?
\end{enumerate}

\begin{enumerate}
\def\labelenumi{(\alph{enumi})}
\tightlist
\item
  Classify each of the six questions using the lecture's taxonomy
  (descriptive, exploratory, inferential, predictive, causal,
  mechanistic). For each, give \textbf{one line} saying what data and
  what kind of analysis would answer it (a summary table, an experiment,
  a model scored on held-out users, \ldots). Name the word or phrase in
  the question that decided your classification.
\end{enumerate}

\begin{enumerate}
\def\labelenumi{(\alph{enumi})}
\setcounter{enumi}{1}
\tightlist
\item
  Question 4 is predictive. Rewrite it as a \textbf{causal} question
  about the same subscription decision --- one sentence. Then say, in
  two or three sentences, what changes: what counts as evidence for the
  predictive version, what counts as evidence for the causal version,
  and why a very accurate churn model does not by itself answer the
  causal one.
\end{enumerate}

\begin{enumerate}
\def\labelenumi{(\alph{enumi})}
\setcounter{enumi}{2}
\tightlist
\item
  A colleague proposes answering question 5 by comparing conversion
  rates of past trial users who \emph{happened} to build their own
  playlist on day 1 against those who did not, and reports ``playlist
  users converted at twice the rate''. Which question type does this
  comparison actually answer? Name one plausible reason the two groups
  differ that has nothing to do with the playlist, and say what design
  would answer question 5 as asked. Keep it to a short paragraph.
\end{enumerate}

\begin{solutionbox}

\begin{enumerate}
\def\labelenumi{(\alph{enumi})}
\tightlist
\item
  One defensible key (the deciding phrase in italics; other readings
  earn credit if argued):
\end{enumerate}

\begin{enumerate}
\def\labelenumi{\arabic{enumi}.}
\tightlist
\item
  \textbf{Descriptive} --- \emph{median \ldots{} last month, by
  country}: a summary of data we already have; a group-by table. No
  claim beyond the users in the table.
\item
  \textbf{Exploratory} --- \emph{natural groups \ldots{} worth a closer
  look}: looking for structure with no labels; a clustering
  (k-means/hierarchical) on genre-listening vectors, followed by plots.
  Generates hypotheses, does not test them.
\item
  \textbf{Inferential} --- \emph{across our whole user base (not just
  the \ldots{} surveyed)}: does a pattern in a sample hold in the
  population? Compare the two groups' podcast rates in the survey with a
  confidence interval or test for the difference.
\item
  \textbf{Predictive} --- \emph{will this user \ldots{} next 30 days}: a
  forecast for a new case; a classifier trained on past users' features
  and churn outcomes, judged by accuracy on held-out users.
\item
  \textbf{Causal} --- \emph{if we give \ldots{} will more of them
  convert}: what happens under an intervention; a randomized experiment
  (A/B test) that assigns the playlist to a random half of trial users.
\item
  \textbf{Mechanistic} --- \emph{by exactly which route}: how the effect
  operates; measure the intermediate steps (session length, daily use)
  and test whether the effect flows through them --- the most demanding
  rung.
\end{enumerate}

\begin{enumerate}
\def\labelenumi{(\alph{enumi})}
\setcounter{enumi}{1}
\item
  Causal rewrite, e.g.~``If we offer \emph{this} user a discount (or a
  reminder, or a new feature) this month, does that reduce the chance
  they cancel?'' --- the key change is an intervention we control.
  Evidence for the predictive version: held-out accuracy --- the model
  is right about who cancels among users it has not seen. Evidence for
  the causal version: that the comparison is free of confounding ---
  ideally an experiment where the intervention is randomly assigned. An
  accurate churn model tells you \emph{which} users are at risk, not
  \emph{what would happen if you acted}: the features it uses (say,
  ``has not opened the app in 10 days'') are correlated with cancelling
  but forcing them to change need not change the outcome. This is the
  lecture's ``type error''.
\item
  The comparison answers an \textbf{inferential/descriptive} association
  question --- ``do users who built playlists convert more?'' --- not
  the causal one asked. Users who build a playlist on day 1 are already
  the engaged ones (more time, more interest in music, more likely to
  pay anyway); that engagement, not the playlist, can explain the
  doubled rate. To answer question 5, randomize: give the playlist to a
  random half of new trial users and withhold it from the other half,
  then compare conversion. Randomization makes the two groups alike in
  everything except the playlist.
\end{enumerate}

\emph{Common mistake:} calling question 3 descriptive because it ``just
compares two averages'' --- the phrase ``across our whole user base''
makes it a claim beyond the sample.

\end{solutionbox}

\answerpage

\subsection{Problem 2: Two vectors by
hand}\label{problem-2-two-vectors-by-hand}

\emph{Practices: norms, dot product, cosine similarity, unit vectors
(A7).}

Let \(\mathbf{x} = (1, 2, 2)\) and \(\mathbf{y} = (2, 2, 1)\) in
\(\mathbb{R}^3\). Definitions you need:
\(\Vert\mathbf{v}\Vert_1 = \sum_i |v_i|\),
\(\Vert\mathbf{v}\Vert_2 = \sqrt{\sum_i v_i^2}\),
\(\Vert\mathbf{v}\Vert_\infty = \max_i |v_i|\),
\(\mathbf{x}\cdot\mathbf{y} = \mathbf{x}^T\mathbf{y} = \sum_i x_i y_i\),
and
\(\cos(\mathbf{x},\mathbf{y}) = \dfrac{\mathbf{x}^T\mathbf{y}}{\Vert\mathbf{x}\Vert_2\,\Vert\mathbf{y}\Vert_2}\).
A \emph{unit vector} has \(\ell_2\) norm 1.

\begin{enumerate}
\def\labelenumi{(\alph{enumi})}
\tightlist
\item
  Compute \(\Vert\mathbf{x}\Vert_1\), \(\Vert\mathbf{x}\Vert_2\),
  \(\Vert\mathbf{x}\Vert_\infty\) and the same three norms of
  \(\mathbf{y}\).
\end{enumerate}

\begin{enumerate}
\def\labelenumi{(\alph{enumi})}
\setcounter{enumi}{1}
\tightlist
\item
  Compute \(\mathbf{x}\cdot\mathbf{y}\), the cosine similarity
  \(\cos(\mathbf{x},\mathbf{y})\) (as an exact fraction), and the angle
  between them (leave it as \(\arccos\) of your fraction, and give
  degrees if you have a calculator). Then write down the unit vector in
  the direction of \(\mathbf{x}\) and verify by direct computation that
  its \(\ell_2\) norm is 1.
\end{enumerate}

\begin{enumerate}
\def\labelenumi{(\alph{enumi})}
\setcounter{enumi}{2}
\tightlist
\item
  Show that for \textbf{every} \(\mathbf{v}\in\mathbb{R}^n\),
  \[\Vert\mathbf{v}\Vert_\infty \;\le\; \Vert\mathbf{v}\Vert_2 \;\le\; \Vert\mathbf{v}\Vert_1 .\]
  A short argument is expected (three or four lines): for the left
  inequality, compare \(\max_i v_i^2\) with \(\sum_i v_i^2\); for the
  right, expand \((\sum_i |v_i|)^2\). Check the chain on your numbers
  from (a).
\end{enumerate}

\begin{enumerate}
\def\labelenumi{(\alph{enumi})}
\setcounter{enumi}{3}
\tightlist
\item
  Give a nonzero vector for which all three norms are \textbf{equal},
  and say in one line what has to be true of a vector for that to
  happen.
\end{enumerate}

\begin{solutionbox}

\begin{enumerate}
\def\labelenumi{(\alph{enumi})}
\item
  \(\Vert\mathbf{x}\Vert_1 = 1+2+2 = 5\);
  \(\Vert\mathbf{x}\Vert_2 = \sqrt{1+4+4} = \sqrt 9 = 3\);
  \(\Vert\mathbf{x}\Vert_\infty = 2\). By symmetry
  \(\Vert\mathbf{y}\Vert_1 = 5\),
  \(\Vert\mathbf{y}\Vert_2 = \sqrt{4+4+1} = 3\),
  \(\Vert\mathbf{y}\Vert_\infty = 2\).
\item
  \(\mathbf{x}\cdot\mathbf{y} = 1\cdot2 + 2\cdot2 + 2\cdot1 = 8\).
  \(\cos(\mathbf{x},\mathbf{y}) = \dfrac{8}{3\cdot 3} = \dfrac{8}{9} \approx 0.889\);
  angle \(\theta = \arccos(8/9) \approx 0.476\) rad
  \(\approx 27.3^\circ\) --- pointing in nearly the same direction. Unit
  vector: \(\mathbf{x}/\Vert\mathbf{x}\Vert_2 = (1/3,\,2/3,\,2/3)\); its
  norm is \(\sqrt{1/9 + 4/9 + 4/9} = \sqrt{9/9} = 1\). (In general
  \(\Vert\mathbf{x}/c\Vert_2 = \Vert\mathbf{x}\Vert_2/c\) for \(c>0\)
  --- the norm's scaling property.)
\item
  Let \(M = \max_i |v_i| = \Vert\mathbf{v}\Vert_\infty\). \emph{Left:}
  \(\Vert\mathbf{v}\Vert_2^2 = \sum_i v_i^2 \ge M^2\) because the sum of
  nonnegative terms is at least its largest term; take square roots.
  \emph{Right:}
  \(\Vert\mathbf{v}\Vert_1^2 = \big(\sum_i |v_i|\big)^2 = \sum_i v_i^2 + \sum_{i\ne j}|v_i||v_j| \ge \sum_i v_i^2 = \Vert\mathbf{v}\Vert_2^2\),
  since the cross terms are nonnegative; take square roots. Check on
  \(\mathbf{x}\): \(2 \le 3 \le 5\). \(\checkmark\)
\item
  Any vector with exactly one nonzero entry,
  e.g.~\(\mathbf{v} = (0,-5,0)\): all three norms equal 5. That is
  exactly when both inequalities in (c) are equalities: one entry
  carries the whole sum (so \(\sum v_i^2 = M^2\)) and every cross term
  \(|v_i||v_j|\), \(i\ne j\), is zero.
\end{enumerate}

\emph{Common mistake:} writing the cosine as \(8/\sqrt{9}\) --- the
denominator is the product of the two norms, \(3\cdot3=9\), not a single
norm.

\end{solutionbox}

\answerpage

\subsection{Problem 3: Data as a
matrix}\label{problem-3-data-as-a-matrix}

\emph{Practices: rows = objects, columns = features; centering; the dot
product of two feature columns as a first measure of ``moving
together''.}

Four students report hours studied and their exam score:

\begin{longtable}[]{@{}ccc@{}}
\toprule\noalign{}
student & hours (\(h\)) & score (\(s\)) \\
\midrule\noalign{}
\endhead
\bottomrule\noalign{}
\endlastfoot
1 & 2 & 60 \\
2 & 4 & 70 \\
3 & 6 & 70 \\
4 & 8 & 80 \\
\end{longtable}

Written as a data matrix \(X\) this is \(4\times 2\): \(m=4\) objects
(rows), \(n=2\) features (columns). To \textbf{mean-center} a column,
subtract that column's mean from each of its entries.

\begin{enumerate}
\def\labelenumi{(\alph{enumi})}
\tightlist
\item
  Compute the mean of each column and write out the centered matrix
  \(X_c\) (the two centered columns \(\mathbf{h}_c\) and
  \(\mathbf{s}_c\)). Verify that each centered column sums to zero, and
  explain in one line why that must always happen.
\end{enumerate}

\begin{enumerate}
\def\labelenumi{(\alph{enumi})}
\setcounter{enumi}{1}
\tightlist
\item
  Compute the dot product \(\mathbf{h}_c\cdot\mathbf{s}_c\) of the two
  centered columns. What does its \textbf{sign} tell you about how the
  two features move together across students? (Think about which
  students contribute positive terms and which negative.)
\end{enumerate}

\begin{enumerate}
\def\labelenumi{(\alph{enumi})}
\setcounter{enumi}{2}
\tightlist
\item
  Compute the cosine similarity of \(\mathbf{h}_c\) and \(\mathbf{s}_c\)
  (three decimals; \(\sqrt{4000} \approx 63.25\)). What would a value
  near \(1\), near \(0\), or near \(-1\) say about hours and scores?
\end{enumerate}

\begin{enumerate}
\def\labelenumi{(\alph{enumi})}
\setcounter{enumi}{3}
\tightlist
\item
  Suppose the scores had been recorded as fractions
  (\(0.60, 0.70, 0.70, 0.80\)) instead of percentages. Without redoing
  everything, say what happens to the dot product in (b) and to the
  cosine in (c), and one sentence on why that makes the cosine the more
  useful of the two as a summary of ``how related''.
\end{enumerate}

\begin{solutionbox}

\begin{enumerate}
\def\labelenumi{(\alph{enumi})}
\item
  \(\bar h = (2+4+6+8)/4 = 5\), \(\bar s = (60+70+70+80)/4 = 70\).
  \[\mathbf{h}_c = (-3,\,-1,\,1,\,3), \qquad \mathbf{s}_c = (-10,\,0,\,0,\,10), \qquad
  X_c = \begin{bmatrix} -3 & -10\\ -1 & 0\\ 1 & 0\\ 3 & 10\end{bmatrix}.\]
  Both columns sum to \(0\):
  \(\sum_i (v_i - \bar v) = \sum_i v_i - m\bar v = m\bar v - m\bar v = 0\)
  by the definition of the mean.
\item
  \(\mathbf{h}_c\cdot\mathbf{s}_c = (-3)(-10) + (-1)(0) + (1)(0) + (3)(10) = 30 + 0 + 0 + 30 = 60\).
  Positive: students below the mean in hours (1) are also below in
  score, and above (4) is above --- the two features move in the same
  direction. A negative product would mean opposite directions; near
  zero, no consistent linear pattern.
\item
  \(\Vert\mathbf{h}_c\Vert_2 = \sqrt{9+1+1+9} = \sqrt{20}\),
  \(\Vert\mathbf{s}_c\Vert_2 = \sqrt{100+0+0+100} = \sqrt{200}\), so
  \(\cos = \dfrac{60}{\sqrt{20}\sqrt{200}} = \dfrac{60}{\sqrt{4000}} \approx \dfrac{60}{63.25} \approx 0.949\).
  Near \(1\): strong same-direction (nearly linear increasing)
  relationship; near \(0\): no linear relationship; near \(-1\): strong
  opposite-direction relationship. (This number \emph{is} the Pearson
  correlation of hours and score --- you will meet it under that name in
  Lecture 05.)
\item
  The centered score column becomes \(\mathbf{s}_c/100\), so the dot
  product shrinks by the same factor to \(60/100 = 0.6\), while the
  cosine is unchanged at \(0.949\): both the numerator and
  \(\Vert\mathbf{s}_c\Vert\) scale by \(1/100\) and the factor cancels.
  The dot product's magnitude depends on the units the features happen
  to be recorded in; the cosine does not --- it measures direction only,
  so it is a unit-free summary of how related the two features are.
\end{enumerate}

\emph{Common mistake:} mean-centering the rows (students) instead of the
columns (features). Centering is per feature.

\end{solutionbox}




\end{document}
